% sprime_lean_section_notes.tex -- what in latex/paper_status.tex (the "Lean Formalization" section) and in
% "figures and tables/architecture.tex" no longer matches the system, now that PA-S' is the source of truth and the
% tau layer lives in it. NOT a rewrite: one item per mismatch, the draft's own words quoted with their line number
% (l. N = line N of latex/paper_status.tex, 158 lines, as of 2026-10-01), and a replacement sentence where needed.
% Notation: $\vdash_k\varphi$ = "the system derives $\varphi$ within budget $k$"; $\models\varphi$ = truth; Lean = prose.
% Fragment: no preamble, plain LaTeX only. Sources: SPRIME_DESIGN_RATIONALE.md, SPRIME_SOURCE_OF_TRUTH.md,
% ROUTE_B_CENSUS.md, MERGE_CROSSROAD.md, app/generated/{sprime_outcome_theorems,sprime_tau_rows}.json, source_diff.md.
\begin{itemize}
\item \textbf{Which system (l.~65, l.~67).} Draft: ``However, this approach requires three costed arithmetic theorems not yet formalized'' and ``we decided to go the other route: we formalize $S$ as an inductive rule set''. The fork is now inverted: the matrix cells and $\tau$-rows that the experiments consume are theorems of the arithmetized system PA-$S'$, and the rule set is at most a comparison. Bounded inner necessitation and the parametric diagonal lemma are theorems and the bounded L\"ob theorem is proved without hypotheses; evaluation facts enter through completeness over the two readings of the action constants, not through a costed $\Sigma_1$-completeness theorem.
\begin{quote}
We formalize $S'$ as Peano Arithmetic with two action constants and a provability predicate bounded by proof length: $\Box_k\sigma$ states that $\sigma$ has a derivation of at most $k$ symbols. The parametric bounded L\"ob theorem is proved for this predicate, for explicit budgets and without hypotheses, following Critch's uniform proof; its one substantial premise, bounded inner necessitation, is itself a theorem, with a polynomial bound of degree 16.
\end{quote}
\item \textbf{The Foundation footnote (l.~65).} Draft: ``but the bound is on the G\"odel number of the derivation rather than on transcript length, and none of the costed derivability conditions are formalized''. The box of $S'$ is built on Foundation's proof predicate and bounds both quantities: the length, and the code, which serves only to make the box a bounded formula and provably excludes no derivation of length at most $k$. The footnote should present Foundation as the library the construction rests on, trusted as a library.
\item \textbf{Figure caption (l.~11).} Draft: ``every matrix cell and $\tau$-row is a theorem, validated by a lint and census inside \texttt{lake build}''. In $S'$ a cell is an outcome theorem or a floor dichotomy (107 and 26 in the library), a $\tau$-row bit is a theorem or a dichotomy bit (149 and 47 of 196), and a third build-time check, coverage, verifies that every pair of the twelve base programs of the paper's zoos is one or the other. What a reader trusts is nine definitions plus the Foundation library; the value used for a floor is a labelled convention applied on the Python side, not a theorem.
\begin{quote}
Architecture and trust boundary. Everything left of the dashed line is checked by the Lean kernel and uses no axiom beyond Lean's three standard ones; a reader trusts the nine definitions of the base and the Foundation library. Every matrix cell is an outcome theorem or a floor dichotomy, and every bit of a $\tau$-row is a theorem or a dichotomy bit; a linter, a census and a coverage check run inside \texttt{lake build}. The two exported JSON files are the only interface: the Python layer performs deterministic arithmetic on them, reads each floor at the outcome of its no-proof branch and labels it, and never reads Lean source.
\end{quote}
\item \textbf{The exported files (l.~11).} Draft: ``The two exported JSON files are the only interface''. Four files now exist: cells and $\tau$-rows of the rule-based system, and cells with floors and $\tau$-rows with floor bits of $S'$. The experiments read only the $S'$ pair; each $\tau$-row file is verified against the cells of its own system when loaded and the two systems are never mixed, so ``two'' is right only if the sentence says which system.
\item \textbf{The architecture figure.} The figure already describes $S'$ (nine definitions, floors as dichotomies, linter, census and coverage) and needs no change of content. Two points do not match the draft: its default section tags (\S4.3 proof methods, \S4.4 floors, \S4.5 $\tau$-rows) assume five subsections where the roadmap sentence (l.~6) announces three, so either the figure's section-tag macros (\texttt{archSecMethods}, \texttt{archSecFloors}, \texttt{archSecTau}) are redefined or the section is restructured; and neither the $\tau$-row box nor the Python side shows that rows contain floor bits read at a labelled branch.
\item \textbf{Third commitment (l.~4).} Draft: ``we use linting to force outcomes onto one typed template with their names matching their content, and a census to verify the converse, that every outcome proven in the library is tagged and linted''. There are now two templates and three checks.
\begin{quote}
Finally, the statements the rest of the paper consumes follow two fixed, enforced templates, one for outcome theorems and one for floor dichotomies: a linter forces every tagged statement onto its template with a name matching its content, a census verifies the converse, and a coverage check verifies that every pair of programs in the paper's zoos is covered by one of the two.
\end{quote}
\item \textbf{The outcome template (l.~37, l.~39--42).} Draft: ``whose first two arguments fix the budget regime (\texttt{.universal}: holds at every budget $k$) and the fuel offset'' and the listing \texttt{OutcomeSpec .universal 2 DupocBot CupodBot (some (.D, .C))}. In $S'$ the statement is about program codes (natural numbers), the evaluator is a relation defined in arithmetic rather than the Lean function \texttt{outcome}, and there is no fuel argument: fuel is existential and the play is unique by determinism.
\begin{quote}
\texttt{@[outcomeA] theorem outcome\_DupocBot\_vs\_CupodBot :}\\
\texttt{\ \ OutcomeSpecA .universal (bot DupocBot) (bot CupodBot) (some (.D, .C))}\\
It unfolds to: for every $k$, the code of \texttt{DupocBot k} run against the code of \texttt{CupodBot k} returns $D$ at some fuel, and conversely $C$.
\end{quote}
\item \textbf{Noncomputability (l.~44).} Draft: ``For that reason, \texttt{ProofSearch} is \emph{noncomputable}, and so is \texttt{outcome}.'' This describes the rule-based evaluator, which calls \texttt{decide} on a proposition. In $S'$ a search node with budget $k$ takes its first branch exactly when $\Box_k$ holds of its guard sentence, and $\Box_k$ quantifies over derivations below an explicit bound.
\begin{quote}
The evaluator is a relation defined in arithmetic, not a Lean function. A guard is a bounded arithmetical fact, but the bound is astronomically large, so outcomes are established by proof and not by running the search.
\end{quote}
\item \textbf{The three layers (l.~50, table caption l.~59).} Draft: ``$\vdash_k\varphi$ is the proposition \texttt{Pf k $\varphi$}, inhabited by an $S$-derivation; $\models\varphi$ is \texttt{$\varphi$.interp}''. In $S'$ layers 0 and 1 are one formula read inside and outside arithmetic; the notation itself is unchanged.
\begin{quote}
$\Box_k\varphi$ is the arithmetical formula ``some derivation of $\varphi$ in the theory has at most $k$ symbols'', which the guard of a program mentions; $\vdash_k\varphi$ states that this formula is true, that is, that such a derivation exists; $\models\varphi$ is the truth of $\varphi$ in $\mathbb{N}$.
\end{quote}
\item \textbf{Rules, soundness per rule, exact costs (l.~71--72, l.~76--109, l.~114).} Draft: ``it is crucial that every rule of $S$ carries its exact transcript cost'' and ``a rule without its case does not compile, so admission into $S$ is the existence of its soundness case''. $S'$ has no rule list (\texttt{Pf}, \texttt{searchBranch}, \texttt{mp}, \texttt{boxIntro} and the \texttt{decide (Pf k \ldots)} clause of l.~85 belong to the rule-based system), one length measure (the symbol count of a derivation) and one soundness theorem. Its results do not depend on exact constants: of the 107 cells, 28 involve no budget, 32 hold at every budget and 47 for all sufficiently large budgets, with thresholds that are existential and never quoted.
\begin{quote}
Soundness is a single theorem: the theory is true in $\mathbb{N}$, so a guard found within its budget states a play that actually occurs. Every negative result, that a guard is not found and the else branch runs, rests on this theorem alone.
\end{quote}
\item \textbf{Faithfulness (l.~115).} Draft: ``which we state as an open obligation'' and ``Any future model, arithmetized or otherwise, inherits every theorem proved here.'' The arithmetized system now exists, but not as a model of the \texttt{BoundedGL} interface at the rule system's costs: there a boxed statement is cited at additive cost, whereas arithmetic pays a polynomial verification, and no budget-preserving translation exists at polynomial budgets. The cells were therefore proved again directly in $S'$, and the outcomes are not all inherited (next item).
\begin{quote}
Faithfulness reduces to nine definitions that a reader checks by eye: the language and theory, the length of a derivation, the box, the coding and description of programs, the guard, substitution, and the evaluator. Everything else is a theorem.
\end{quote}
\item \textbf{Cells on which the two systems differ (no sentence in the draft; contradicts l.~115).} Of the 107 cells stated in $S'$, 94 agree with the rule-based system and 13 contradict it: 11 involve \texttt{DIMCID} and 2 are \texttt{WaryBot}'s (against \texttt{MirrorBot} and against itself). Inside the paper's zoos there are 9, all \texttt{DIMCID}'s (7 on the body zoo), and among the 196 $\tau$-row bits there are 13, all involving \texttt{DIMCID}. The cause is that arithmetic reads every source, while the rule-based proofs of these cells rest on a program being unreadable by the rules.
\begin{quote}
In $S'$, \texttt{DIMCID} defects against every closed program: its guard, an implication whose antecedent is its own cooperation, is self-fulfilling, and the bounded L\"ob theorem makes it provable within the budget for all sufficiently large $k$.
\end{quote}
\item \textbf{The pricing decision (l.~118, l.~133).} Draft: ``In our case, we chose to price every failure certificate strictly above the failed budget $k$'' and ``Whether a finer pricing is consistent is left open.'' $S'$ has nothing to price; the anti-diagonal argument (l.~119--131) stays valid as the reason why no system certifies all of its own failed searches at the same budget.
\begin{quote}
For 26 cells the outcome hinges on whether a true guard sentence has a proof within the budget $k$. For each we prove a dichotomy without hypotheses (not found: one outcome; found: another), and we show on a representative cell that the hinge is equivalent, up to an additive term, to the consistency of the theory for proofs of length $k$. Pudl\'ak's lower bound does not settle this regime, so neither branch is assumed; the experiments use the not-found branch and label the cell.
\end{quote}
\item \textbf{\texttt{DupocBot} against \texttt{CupodTrollBot}, and staggering (l.~135, l.~143).} Draft: ``\texttt{DupocBot} does not fire its cooperation guard and defects'' and ``\texttt{DupocBot} $k$ defects against \texttt{CupodTrollBot} $k$, whereas \texttt{DupocBot} $(2k {+} 64)$ cooperates, for large $k$.'' In $S'$ this cell is one of the 26 floors, and the staggered statement is a theorem of the rule-based system only (the $S'$ export contains no staggered companion).
\begin{quote}
Against \texttt{CupodTrollBot}, the guard of \texttt{DupocBot} is true, since \texttt{CupodTrollBot} cooperates at every budget. Whether that guard has a proof within $k$ is the hinge: if not, the outcome is $(D,C)$; if so, it is $(C,C)$.
\end{quote}
\item \textbf{\texttt{PrudentBot} (l.~156, l.~158).} Draft: ``\texttt{PrudentBot} $k$ defects against itself, but \texttt{PrudentBot2}$(k,4k{+}100)$ cooperates against itself, for large $k$'' and ``With our rule, self-fulfilling facts are certified at one tier, but a failure fact can only be certified by a strictly higher tier, about a lower one.'' In $S'$ the self-play of \texttt{PrudentBot} is a floor dichotomy, $(D,D)$ if its guard is not found within $k$ and $(C,C)$ if it is; the two-tier theorem and ``our rule'' are statements about the rule-based system. The analogy with PA${+}1$ survives as a question rather than a rule: how long is the shortest proof that a search of budget $k$ fails.
\item \textbf{A program inside a guard (l.~137--141, \texttt{.eq .opp (CupodBot k)}).} The draft does not say how a guard describes a program that carries the budget, and in $S'$ this is part of the trusted base. Written out literally, such a description contains the numeral of a formula that changes with $k$, of length at least $k$, and the guard could never be found.
\begin{quote}
A program that occurs inside a guard is described by a term; where that program carries the guard's own budget, the description names the budget by a variable, which is instantiated when the guard is consulted. Instantiating the variable gives back the literal description, so the sentence consulted is the same; the gain is that the stored guard is one fixed formula at every budget, with a description of length logarithmic in $k$.
\end{quote}
\item \textbf{$\tau$-rows (l.~6, l.~11).} The partial-transparency subsection announced at l.~6 is not in the draft; when written, its rows are statements of $S'$. For the 14 templates there are 196 bits: 149 are theorems and 47 are floor bits, six of which are images of the Dupoc--Cupod cell, a theorem in the base matrix but a dichotomy on the compiled two-member systems. The check that the lifted programs agree with the base cells gives 149 agreements, 41 agreements on floor bits, those 6 bits with the same value but a different status, and no mismatch.
\item \textbf{Open obligations (footnote, l.~4).} Draft: ``Some meta-level obligations are left open''. The list in the Conclusion changes: an arithmetical model of the interface is no longer the open item. What stays open is the hinge of the 26 floors (one question about proof length), the programs with a box inside a guard (\texttt{LegibleBot}, \texttt{OptimBot}), which are outside the coverage of $S'$, and the trust placed in Foundation.
\item \textbf{Checked, no edit needed.} ``we do not introduce any axioms, apart from Lean's three standard axioms'' and the \texttt{sorry}-free claim (l.~4) hold for $S'$. ``\texttt{Prog} has ten constructors in total'' (footnote, l.~19) is still true of the surface syntax, and program codes have the ten matching shapes; of the three partial-transparency constructors only two occur in the programs the $\tau$-rows are about. The listings of \texttt{DupocBot}, \texttt{DBot}, \texttt{CupodTrollBot} and \texttt{PrudentBot2} are unchanged.
\end{itemize}
